% 付録・第1部：行列の基本から漸化式へ

\section{行列の基本から,特性方程式を見直す}\label{節：特性方程式を行列で見直す}

本節では,まず行列の読み方と計算の仕方を説明する.
その後,本編で使った「等比数列になるように置き換える」という工夫を,行列の言葉で見直していこう.
行列を初めて学ぶ読者も,最初の数値例から順に確かめてほしい.
ここでは主に二行二列の行列を扱う.

\subsection{行列とは何か}

行列とは,数を長方形に並べ,全体を括弧で囲んだものである.
たとえば
\[
  A=\begin{pmatrix}2&1\\1&2\end{pmatrix}
\]
は,縦に二段,横に二列の数を並べた行列である.
横の並びを\textbf{行},縦の並びを\textbf{列}という.
上が第一行,下が第二行であり,左が第一列,右が第二列である.
並んでいる一つ一つの数を\textbf{成分}と呼ぶ.
この例の第一行・第二列の成分は$1$である.

また,二つの数を縦に並べた
\[
  \boldsymbol{v}=\begin{pmatrix}x\\y\end{pmatrix}
\]
を\textbf{列ベクトル}と呼ぶ.
太字の$\boldsymbol{v}$は,二つの数をまとめて表すための記号である.
平面上の点$(x,y)$への矢印と考えることもできるが,以下ではまず「二つの値を入れた一組の記録」として読めばよい.

列ベクトルの足し算と定数倍は,成分ごとに行う.
\[
  \begin{pmatrix}x\\y\end{pmatrix}
  +\begin{pmatrix}u\\v\end{pmatrix}
  =\begin{pmatrix}x+u\\y+v\end{pmatrix},
  \qquad
  c\begin{pmatrix}x\\y\end{pmatrix}
  =\begin{pmatrix}cx\\cy\end{pmatrix}.
\]
たとえば$\begin{pmatrix}3\\1\end{pmatrix}+\begin{pmatrix}1\\-1\end{pmatrix}
=\begin{pmatrix}4\\0\end{pmatrix}$である.
二つの列ベクトルが等しいとは,上下それぞれの成分が等しいことをいう.

\subsubsection{行列を列ベクトルに掛ける}

先ほどの行列$A$を列ベクトルに掛けるときは,次のように計算する.
\[
  \begin{pmatrix}2&1\\1&2\end{pmatrix}
  \begin{pmatrix}x\\y\end{pmatrix}
  =
  \begin{pmatrix}2x+y\\x+2y\end{pmatrix}.
\]
上の成分は,第一行の$2,\,1$を$x,y$にそれぞれ掛けて足したもの.
下の成分は,第二行の$1,\,2$を$x,y$にそれぞれ掛けて足したものである.
たとえば
\[
  A\begin{pmatrix}1\\0\end{pmatrix}
  =\begin{pmatrix}2\\1\end{pmatrix},
  \qquad
  A\begin{pmatrix}2\\1\end{pmatrix}
  =\begin{pmatrix}5\\4\end{pmatrix}.
\]
この行列は,二つの値$x,y$を受け取り,新しい二つの値$2x+y,x+2y$を作る規則を表している.

一般には
\[
  \begin{pmatrix}p&q\\r&s\end{pmatrix}
  \begin{pmatrix}x\\y\end{pmatrix}
  =\begin{pmatrix}px+qy\\rx+sy\end{pmatrix}
\]
と定める.
したがって,連立した二つの式
\[
  x'=px+qy,\qquad y'=rx+sy
\]
を,行列と列ベクトルの積一つにまとめられる.
ここで$x',y'$は,更新後の値を表す記号である.

また,成分ごとに展開すれば
\[
  A(c\boldsymbol{u}+d\boldsymbol{v})
  =cA\boldsymbol{u}+dA\boldsymbol{v}
\]
が成り立つ.
これは,二つの列ベクトルを足してから更新しても,別々に更新してから足しても結果が同じ,という性質である.
後で初期値を二つの部分に分けて考える際に使う.

\subsection{行列の積は,更新を続けて行うことを表す}

行列をもう一度掛けると,同じ規則で二回更新できる.
先ほどの例なら
\[
  \begin{pmatrix}1\\0\end{pmatrix}
  \ \longmapsto\
  \begin{pmatrix}2\\1\end{pmatrix}
  \ \longmapsto\
  \begin{pmatrix}5\\4\end{pmatrix}
\]
と進む.
この二回分の更新を一つにまとめた行列を$A^2$と書く.
実際に文字$x,y$で計算すると
\[
  \begin{aligned}
    A\left(A\begin{pmatrix}x\\y\end{pmatrix}\right)
    &=A\begin{pmatrix}2x+y\\x+2y\end{pmatrix}\\
    &=\begin{pmatrix}5x+4y\\4x+5y\end{pmatrix}.
  \end{aligned}
\]
したがって
\[
  A^2=\begin{pmatrix}5&4\\4&5\end{pmatrix}
\]
である.
各成分を単に二乗するのではなく,更新を二回行った結果に合わせて積を定めている.

異なる行列$A,B$でも,まず$B$,次に$A$を掛ける操作を$AB$と書く.
$B$の各列に$A$を掛けた結果を並べると,積$AB$が得られる.
具体的には
\[
  \begin{aligned}
    &\begin{pmatrix}p&q\\r&s\end{pmatrix}
     \begin{pmatrix}u&v\\w&z\end{pmatrix}\\
    &\qquad=
     \begin{pmatrix}
       pu+qw&pv+qz\\
       ru+sw&rv+sz
     \end{pmatrix}.
  \end{aligned}
\]
右の行列から先に作用するため,一般には$AB$と$BA$は一致しない.

何も変えない更新は
\[
  I=\begin{pmatrix}1&0\\0&1\end{pmatrix},
  \qquad
  I\begin{pmatrix}x\\y\end{pmatrix}
  =\begin{pmatrix}x\\y\end{pmatrix}
\]
で表される.
$I$を\textbf{単位行列}という.
数の累乗で$a^0=1$とするのに対応して,行列では$A^0=I$と定める.

ここで,二つの数列が
\[
  x_{n+1}=2x_n+y_n,\qquad
  y_{n+1}=x_n+2y_n
\]
を満たす場合を考えよう.
$n$番目の二つの値を
\[
  \boldsymbol{v}_n=\begin{pmatrix}x_n\\y_n\end{pmatrix}
\]
とまとめれば,\,$\boldsymbol{v}_{n+1}=A\boldsymbol{v}_n$である.
このように,次の値を計算するために持っておく値の組を\textbf{状態}と呼ぶ.
初期状態から一回,二回と進めると
\[
  \boldsymbol{v}_2=A\boldsymbol{v}_1,\qquad
  \boldsymbol{v}_3=A^2\boldsymbol{v}_1
\]
なので,一般に
\[
  \boldsymbol{v}_n=A^{n-1}\boldsymbol{v}_1
\]
となる.
初項から第$n$項までの更新は$n-1$回である.
行列の累乗が分かれば,連立漸化式の一般項も分かることになる.

\subsection{固有ベクトルと固有値：更新が定数倍になる場合}

毎回行列の積を計算するのは手間がかかる.
そこで,行列を掛けても定数倍になるだけの列ベクトルを探してみよう.
先ほどの$A$について
\[
  \begin{aligned}
    A\begin{pmatrix}1\\1\end{pmatrix}
      &=\begin{pmatrix}3\\3\end{pmatrix}
       =3\begin{pmatrix}1\\1\end{pmatrix},\\
    A\begin{pmatrix}1\\-1\end{pmatrix}
      &=\begin{pmatrix}1\\-1\end{pmatrix}
  \end{aligned}
\]
である.
第一の列ベクトルは一回ごとに$3$倍され,第二の列ベクトルは変わらない.
したがって$k$回更新した結果は
\[
  A^k\begin{pmatrix}1\\1\end{pmatrix}
    =3^k\begin{pmatrix}1\\1\end{pmatrix},
  \qquad
  A^k\begin{pmatrix}1\\-1\end{pmatrix}
    =\begin{pmatrix}1\\-1\end{pmatrix}
\]
とすぐに分かる.

一般に,零でない列ベクトル$\boldsymbol{u}$が
\[
  A\boldsymbol{u}=\lambda\boldsymbol{u}
\]
を満たすとき,\,$\boldsymbol{u}$を$A$の\textbf{固有ベクトル},\,$\lambda$を対応する\textbf{固有値}という.
固有値は,そのベクトルに対して一回の更新が何倍になるかを表す数である.
零ベクトル$\boldsymbol{0}=\begin{pmatrix}0\\0\end{pmatrix}$は,どんな$\lambda$でもこの式を満たしてしまうため,固有ベクトルから除く.
負の固有値では矢印の向きが反転し,固有値$0$では零ベクトルになる場合もある.

\subsubsection{初期状態を二つに分けると,一般項が出る}

初期値を$(x_1,y_1)=(1,\,0)$として,先ほどの連立漸化式を解いてみよう.
初期状態は
\[
  \begin{pmatrix}1\\0\end{pmatrix}
  =\frac12\begin{pmatrix}1\\1\end{pmatrix}
   +\frac12\begin{pmatrix}1\\-1\end{pmatrix}
\]
と分けられる.
右辺を実際に足せば,上が$1$,下が$0$になる.
二つの部分を別々に$n-1$回更新すると
\[
  \begin{pmatrix}x_n\\y_n\end{pmatrix}
  =\frac{3^{n-1}}2\begin{pmatrix}1\\1\end{pmatrix}
   +\frac12\begin{pmatrix}1\\-1\end{pmatrix}.
\]
上下の成分を読み取って
\[
  x_n=\frac{3^{n-1}+1}{2},\qquad
  y_n=\frac{3^{n-1}-1}{2}
\]
を得る.
初期状態を,更新の簡単な二つの部分に分けたことが解法の要点である.

本編の連立型では,和と差を
\[
  u_n=x_n+y_n,\qquad v_n=x_n-y_n
\]
と置いた.
もとの二つの式を足す,引くという計算から
\[
  u_{n+1}=3u_n,\qquad v_{n+1}=v_n
\]
となる.
そして
\[
  x_n=\frac{u_n+v_n}{2},\qquad
  y_n=\frac{u_n-v_n}{2}
\]
で元に戻せる.
ここに現れた倍率$3,\,1$は,先ほどの固有値と同じである.

和と差を作る係数を横に並べると,\,$(1,\,1)$と$(1,-1)$になる.
このように数を横に並べたものを\textbf{行ベクトル}という.
行ベクトルと列ベクトルの積は,対応する成分を掛けて足す.
\[
  (\alpha,\beta)\begin{pmatrix}x\\y\end{pmatrix}
  =\alpha x+\beta y.
\]
行ベクトルを行列に掛けるときは,行列の各列との積を横に並べる.
たとえば$(1,\,1)A=(3,\,3)=3(1,\,1)$である.
一般の行列$A$でも,零でない行ベクトル$L$が
\[
  LA=\gamma L
\]
を満たせば,\,$c_n=L\boldsymbol{v}_n$は
\[
  c_{n+1}=LA\boldsymbol{v}_n=\gamma c_n
\]
を満たす.
この$L$を\textbf{左固有ベクトル}と呼ぶ.
本編で等比数列になるように$\alpha x_n+\beta y_n$の係数を選ぶ方法は,この行ベクトルを探す計算として理解できる.
固有ベクトルで状態を分ける方法と,等比数列になる量を取り出す方法は,同じ更新を異なる側から見ている.

\subsection{本編の特性方程式型を,行列で書き直す}\label{節：一つの値の更新を行列で書く}

次に,本編で扱った
\[
  a_{n+1}=pa_n+q\qquad(p\ne1)
\]
に戻ろう.
ここでは一つの値だけを更新しているが,定数$q$の足し算がある.
この足し算を行列の掛け算の中に含めるため,\,$a_n$と一緒に$1$も並べておく.
すると
\[
  \begin{pmatrix}p&q\\0&1\end{pmatrix}
  \begin{pmatrix}a_n\\1\end{pmatrix}
  =\begin{pmatrix}pa_n+q\\1\end{pmatrix}
  =\begin{pmatrix}a_{n+1}\\1\end{pmatrix}.
\]
上の成分では$q$に$1$を掛けることで,必要な足し算を作っている.
下の成分は毎回$1$のままなので,次の更新でも同じ計算ができる.

この行列を$H$とし,本編と同じく
\[
  \alpha=p\alpha+q,\qquad
  \alpha=\frac{q}{1-p}
\]
を満たす$\alpha$を取る.
この値から始めると次の項も変わらないので,\,$\alpha$は不動点である.
行列で計算すると
\[
  H\begin{pmatrix}\alpha\\1\end{pmatrix}
  =\begin{pmatrix}\alpha\\1\end{pmatrix},
  \qquad
  H\begin{pmatrix}1\\0\end{pmatrix}
  =p\begin{pmatrix}1\\0\end{pmatrix}.
\]
これで,更新が簡単な二つの列ベクトルが見つかった.

初期状態を
\[
  \begin{pmatrix}a_1\\1\end{pmatrix}
  =\begin{pmatrix}\alpha\\1\end{pmatrix}
   +(a_1-\alpha)\begin{pmatrix}1\\0\end{pmatrix}
\]
と分ける.
第一の部分は変わらず,第二の部分だけが一回ごとに$p$倍されるから
\[
  \begin{pmatrix}a_n\\1\end{pmatrix}
  =\begin{pmatrix}\alpha\\1\end{pmatrix}
   +(a_1-\alpha)p^{n-1}\begin{pmatrix}1\\0\end{pmatrix}.
\]
上の成分を読めば
\[
  a_n=\alpha+(a_1-\alpha)p^{n-1}
\]
となる.
$p=0$では$a_2$以降が$\alpha=q$になるので直接読めばよい.

たとえば$a_{n+1}=2a_n+1,\ a_1=0$なら$\alpha=-1$である.
初期状態は
\[
  \begin{pmatrix}0\\1\end{pmatrix}
  =\begin{pmatrix}-1\\1\end{pmatrix}
   +\begin{pmatrix}1\\0\end{pmatrix}
\]
と分かれ,後半だけが$2$倍ずつになる.
したがって$a_n=-1+2^{n-1}$である.
本編の$b_n=a_n-\alpha$という置き換えは,この「不動点からのずれ」だけを取り出していたのである.

ここで用いた$\alpha=p\alpha+q$は,不動点を求める方程式である.
一方,行列の固有値を求める方程式も線形代数では「特性方程式」と呼ばれる.
二つは異なる式であり,不動点$\alpha$そのものを固有値と呼ぶわけではない.
上の行列$H$の固有値は$1$と$p$である.

$p=1$なら,もとの式は$a_{n+1}=a_n+q$である.
更新を$k$回続けると
\[
  H^k=\begin{pmatrix}1&kq\\0&1\end{pmatrix}
\]
となり,\,$a_n=a_1+(n-1)q$を得る.
$q\ne0$では不動点がなく,加えた$q$が更新回数だけ積み重なる.
これは後で見る,特性根が重なると一般項に$n$が現れる例につながっている.

\subsection{二つの項を記録すれば,三項間漸化式も行列になる}\label{節：状態を二つ持つと,二階の漸化式になる}

三項間漸化式では,次の項を計算するために直前の二つの項が必要である.
たとえばフィボナッチ数列
\[
  F_0=0,\quad F_1=1,\quad
  F_{n+2}=F_{n+1}+F_n\quad(n\ge0)
\]
を考える.
状態を$\begin{pmatrix}F_{n+1}\\F_n\end{pmatrix}$とすれば
\[
  \begin{pmatrix}1&1\\1&0\end{pmatrix}
  \begin{pmatrix}F_{n+1}\\F_n\end{pmatrix}
  =\begin{pmatrix}F_{n+1}+F_n\\F_{n+1}\end{pmatrix}
  =\begin{pmatrix}F_{n+2}\\F_{n+1}\end{pmatrix}.
\]
上の成分で新しい項を作り,下の成分に直前の項を残している.
実際,
\[
  \begin{pmatrix}1\\0\end{pmatrix}
  \longmapsto\begin{pmatrix}1\\1\end{pmatrix}
  \longmapsto\begin{pmatrix}2\\1\end{pmatrix}
  \longmapsto\begin{pmatrix}3\\2\end{pmatrix}
\]
と進む.

一般に
\[
  a_{n+2}=s a_{n+1}+t a_n\qquad(n\ge0)
\]
は
\[
  \boldsymbol{w}_n=\begin{pmatrix}a_{n+1}\\a_n\end{pmatrix},
  \qquad
  B=\begin{pmatrix}s&t\\1&0\end{pmatrix}
\]
と置けば,\,$\boldsymbol{w}_{n+1}=B\boldsymbol{w}_n$になる.
直前の二項を必要とするため,この形を\textbf{二階の漸化式}とも呼ぶ.
本編の$a_{n+2}+pa_{n+1}+qa_n=0$では,\,$s=-p,\ t=-q$とすればよい.

\subsubsection{なぜ本編と同じ特性方程式が出るのか}

$B$を掛けると$\lambda$倍になる,零でない列ベクトルを探そう.
成分を$x,y$と書けば
\[
  B\begin{pmatrix}x\\y\end{pmatrix}
  =\lambda\begin{pmatrix}x\\y\end{pmatrix}
  \quad\Longleftrightarrow\quad
  \begin{cases}
    sx+ty=\lambda x,\\
    x=\lambda y.
  \end{cases}
\]
二つ目の式から,\,$y=0$なら$x=0$になってしまう.
そこで$y\ne0$と分かるので,ベクトル全体を$y$で割り,\,$y=1,\ x=\lambda$としてよい.
一つ目の式は
\[
  s\lambda+t=\lambda^2,
  \qquad
  \lambda^2-s\lambda-t=0
\]
となる.
これが固有値を求める方程式である.

本編では,等比数列の形$a_n=\lambda^n$を漸化式に当てはめて,同じ方程式を得た.
行列では「二項をまとめた記録が一回ごとに$\lambda$倍になる条件」を調べている.
数列が等比的に進むことと,状態ベクトルが定数倍になることは,同じ条件を別の形で表している.
$\lambda=0$が根になる場合も,行列の方程式はそのまま使える.

二つの異なる根$\lambda_1,\lambda_2$があれば
\[
  \begin{pmatrix}\lambda_1\\1\end{pmatrix},
  \qquad
  \begin{pmatrix}\lambda_2\\1\end{pmatrix}
\]
が固有ベクトルになる.
初期状態を
\[
  \begin{pmatrix}a_1\\a_0\end{pmatrix}
  =c\begin{pmatrix}\lambda_1\\1\end{pmatrix}
   +d\begin{pmatrix}\lambda_2\\1\end{pmatrix}
\]
と分けるには,\,$c+d=a_0,\ c\lambda_1+d\lambda_2=a_1$を解けばよい.
根が異なるので,\,$c,d$は一意に決まる.
$n$回更新した後は,下の成分から
\[
  a_n=c\lambda_1^n+d\lambda_2^n
\]
を得る.
根が$0$の場合,\,$n=0$での係数は初期状態の式から読み,\,$n\ge1$ではその成分は$0$になる.

フィボナッチ数列なら特性方程式は$\lambda^2-\lambda-1=0$で,根は
\[
  \varphi=\frac{1+\sqrt5}{2},\qquad
  \psi=\frac{1-\sqrt5}{2}
\]
である.
初期状態$\begin{pmatrix}1\\0\end{pmatrix}$に合わせる条件は
\[
  c+d=0,\qquad c\varphi+d\psi=1.
\]
$\varphi-\psi=\sqrt5$より$c=1/\sqrt5,\ d=-1/\sqrt5$と決まり,
\[
  F_n=\frac{\varphi^n-\psi^n}{\sqrt5}
\]
を得る.
$|\psi|<1<\varphi$なので,\,$\psi^n$の成分は小さくなり,\,$\varphi^n$の成分が大きくなる.
固有値は一般項の計算に加えて,更新を重ねたときの増え方も教えている.
異なる二つの根が複素数の場合も,成分や係数を複素数まで広げれば同じ計算ができる.
実数の初期値に合わせると,共役な二つの根の寄与が組になり,数列の項は実数になる.

\subsubsection{行列式を使うと,固有値の条件をまとめられる}

ここまでのように成分を比較すれば,固有値は求められる.
その条件を一つの式にまとめる記号が\textbf{行列式}である.
二行二列の場合には
\[
  \det\begin{pmatrix}a&b\\c&d\end{pmatrix}=ad-bc
\]
と定める.
行列そのものではなく,四つの成分から計算した一つの数である.

二つの式$ax+by=0,\ cx+dy=0$について,\,$ad-bc\ne0$なら解は$x=y=0$だけになる.
たとえば,第一の式を$d$倍したものから第二の式を$b$倍したものを引くと,\,$(ad-bc)x=0$となる.
同様に$(ad-bc)y=0$も得られる.
逆に$ad-bc=0$なら,零でない解が存在する.
実際,\,$(a,b)\ne(0,\,0)$なら$(x,y)=(b,-a)$が両方の式を満たす.
$(a,b)=(0,\,0)$の場合も,残る一つの式から零でない解を選べる.

固有値の式$B\boldsymbol{u}=\lambda\boldsymbol{u}$は,
\[
  (\lambda I-B)\boldsymbol{u}=\boldsymbol{0}
\]
と書き直せる.
行列の引き算は対応する成分ごとに行うので
\[
  \lambda I-B
  =\begin{pmatrix}\lambda-s&-t\\-1&\lambda\end{pmatrix}.
\]
零でない解が存在する条件は
\[
  \begin{aligned}
    \det(\lambda I-B)
      &=(\lambda-s)\lambda-(-t)(-1)\\
      &=\lambda^2-s\lambda-t=0.
  \end{aligned}
\]
こうして,先ほど成分から求めた特性方程式が再び得られる.
一般の行列$A=\begin{pmatrix}p&q\\r&s\end{pmatrix}$でも,同じ計算で
\[
  \det(\lambda I-A)
  =\lambda^2-(p+s)\lambda+(ps-qr)=0
\]
が固有値を求める特性方程式になる.

\subsubsection{特性根が重なるとき}

特性方程式が重解$r\ne0$を持つときには,異なる二つの固有ベクトルに分ける先ほどの方法が使えない.
一般項に$n$が現れる理由は,もとの漸化式を直接変形すると分かりやすい.
重解のときは
\[
  a_{n+2}=2r a_{n+1}-r^2a_n
\]
なので,\,$a_n=r^n b_n$と置いて$r^{n+2}$で割ると
\[
  b_{n+2}-2b_{n+1}+b_n=0.
\]
これは
\[
  b_{n+2}-b_{n+1}=b_{n+1}-b_n
\]
ということだから,\,$b_n$は等差数列になる.
したがって
\[
  b_n=C+Dn,\qquad a_n=(C+Dn)r^n
\]
である.
等比的な倍率$r^n$を取り除くと,残りに等差的な増え方が現れるのである.
$r=0$が重解なら$a_{n+2}=0$なので,\,$a_2$以降はすべて$0$として直接扱う.

\subsection{一次分数型は,二つの成分の比として読む}

ここまで,列ベクトルの成分そのものを数列の項として読んだ.
今度は成分の\textbf{比}を読むと,本編の一次分数型につながることを見てみよう.
\[
  a_{n+1}=\frac{pa_n+q}{ra_n+s}
\]
に対し,\,$r\ne0,\ ps-qr\ne0$とし,すべての項で$ra_n+s\ne0$と仮定する.
初めを$(x_1,y_1)=(a_1,\,1)$として
\[
  \begin{pmatrix}x_{n+1}\\y_{n+1}\end{pmatrix}
  =M\begin{pmatrix}x_n\\y_n\end{pmatrix},
  \qquad M=\begin{pmatrix}p&q\\r&s\end{pmatrix}
\]
と更新する.
$a_n=x_n/y_n$であれば
\[
  \begin{aligned}
    \frac{x_{n+1}}{y_{n+1}}
    &=\frac{px_n+qy_n}{rx_n+sy_n}\\
    &=\frac{p(x_n/y_n)+q}{r(x_n/y_n)+s}
     =\frac{pa_n+q}{ra_n+s}.
  \end{aligned}
\]
また,\,$y_{n+1}=(ra_n+s)y_n$なので,分母が零にならないという仮定のもとで$y_n\ne0$が保たれる.
したがって,各回の成分比はもとの漸化式の項になる.
ベクトル全体を零でない数で何倍しても比は同じである.
このように成分比に注目して行列の作用を見る方法を,射影変換の見方という.

不動点$\alpha$については
\[
  p\alpha+q=\alpha(r\alpha+s)
\]
なので
\[
  M\begin{pmatrix}\alpha\\1\end{pmatrix}
  =(r\alpha+s)\begin{pmatrix}\alpha\\1\end{pmatrix}.
\]
つまり,不動点に対応する列ベクトルは固有ベクトルであり,固有値は$r\alpha+s$である.
$\alpha$自身が固有値なのではなく,ベクトルの二成分の比になっている.

異なる二つの不動点$\alpha,\beta$がある場合を考え,
\[
  \lambda_\alpha=r\alpha+s,\qquad
  \lambda_\beta=r\beta+s
\]
と置く.
各状態を
\[
  \begin{pmatrix}x_n\\y_n\end{pmatrix}
  =c_n\begin{pmatrix}\alpha\\1\end{pmatrix}
   +d_n\begin{pmatrix}\beta\\1\end{pmatrix}
\]
と分ければ
\[
  c_{n+1}=\lambda_\alpha c_n,\qquad
  d_{n+1}=\lambda_\beta d_n
\]
となる.
二つの係数が,別々の等比数列として進むのである.

本編の置き換えとの関係を,成分から確かめよう.
$x_n=\alpha c_n+\beta d_n,\ y_n=c_n+d_n$だから
\[
  x_n-\alpha y_n=(\beta-\alpha)d_n,
  \qquad
  x_n-\beta y_n=(\alpha-\beta)c_n.
\]
したがって
\[
  \frac{a_n-\alpha}{a_n-\beta}
  =\frac{x_n-\alpha y_n}{x_n-\beta y_n}
  =-\frac{d_n}{c_n}.
\]
以下では$a_1\ne\beta$とする.
このとき$c_1\ne0$であり,\,$ps-qr\ne0$より$\lambda_\alpha\ne0$なので,\,$c_n\ne0$が保たれる.
ここで$\lambda_\alpha=0$なら$M\begin{pmatrix}\alpha\\1\end{pmatrix}=\boldsymbol{0}$となり,行列式が零でないことに反する.
$a_1=\beta$の場合は,\,$a_n=\beta$という定数列として扱えばよい.

係数の更新規則より
\[
  \frac{a_{n+1}-\alpha}{a_{n+1}-\beta}
  =\frac{\lambda_\beta}{\lambda_\alpha}
   \frac{a_n-\alpha}{a_n-\beta}.
\]
本編で$b_n=(a_n-\alpha)/(a_n-\beta)$を等比数列にできたのは,この二つの係数の比を取り出していたからである.
不動点の方程式
\[
  r z^2+(s-p)z-q=0
\]
の解と係数の関係から$r(\alpha+\beta)=p-s$なので,
\[
  p-r\alpha=r\beta+s,\qquad
  p-r\beta=r\alpha+s
\]
となる.
よって公比は本編と同じ
\[
  \frac{\lambda_\beta}{\lambda_\alpha}
  =\frac{p-r\alpha}{p-r\beta}
\]
である.

たとえば$a_{n+1}=2a_n/(1-a_n)$では
\[
  M=\begin{pmatrix}2&0\\-1&1\end{pmatrix}
\]
で,不動点$0,-1$に対応する固有値はそれぞれ$1,\,2$になる.
$\alpha=0,\ \beta=-1$とすると
\[
  b_n=\frac{a_n}{a_n+1},\qquad b_{n+1}=2b_n
\]
である.
この$2$は,二つの固有値の比$2/1$から現れていた.

\subsection{三角関数型漸化式からチェビシェフ多項式へ}\label{節：三角関数型漸化式からチェビシェフ多項式へ}

特性根が複素数になると,数列には振動が現れる.
実係数の漸化式では複素根は共役な組で現れ,一般には$\rho>0,\ 0<\theta<\pi$の条件下で
\[
  \rho(\cos\theta+i\sin\theta),
  \qquad
  \rho(\cos\theta-i\sin\theta)
\]
と書ける.
特性方程式の係数を比べると,対応する漸化式は
\[
  a_{n+2}=2\rho\cos\theta\,a_{n+1}-\rho^2a_n
\]
である.
ド・モアブルの公式から,特性根の $n$ 乗は振幅 $\rho^n$ と角度 $n\theta$ を持つ.
そのため実数解は
\[
  a_n=\rho^n\{C\cos(n\theta)+D\sin(n\theta)\}
\]
という形になる.
$\rho>1$ なら振幅は大きくなり,$\,0<\rho<1$ なら小さくなる.

特に振幅が変わらない $\rho=1$ の場合,特性根は $\cos\theta\pm i\sin\theta$ であり,漸化式は
\[
  a_{n+2}=2\cos\theta\,a_{n+1}-a_n
\]
となる.
余弦と正弦の加法定理を使うと,$\,\cos(n\theta)$ も $\,\sin(n\theta)$ もこの同じ漸化式を満たす.
したがって,二つの初期値に合わせて定数 $C,D$ を選べば
\[
  a_n=C\cos(n\theta)+D\sin(n\theta)
\]
と表せる.
複素数の根は,ただ計算を複雑にするものではない.
実数の数列にひそむ振動と,その振幅の増減を記録している.

ここで $x=\cos\theta$ と書き,角度の整数倍を $x$ の多項式で表せないか考えてみよう.
余弦の加法定理
\[
  \cos((n+1)\theta)+\cos((n-1)\theta)
  =2\cos\theta\cos(n\theta)
\]
を変形すると
\[
  \cos((n+1)\theta)
  =2x\cos(n\theta)-\cos((n-1)\theta)
\]
となる.
右辺は,一つ前と二つ前の余弦から次の余弦を作っている.
そこで同じ規則を多項式の列に移し,
\begin{align*}
  &T_0(x)=1,\qquad T_1(x)=x,\\
  &T_{n+1}(x)=2xT_n(x)-T_{n-1}(x)
  \qquad(n\ge1)
\end{align*}

と定める.
最初の二つが多項式で,次も足し算と掛け算だけで作られるので,すべての $T_n(x)$ は多項式になる.
この列をチェビシェフ多項式という.

この定義から
\[
  T_2(x)=2x^2-1,\qquad
  T_3(x)=4x^3-3x
\]
が得られる.
これらは二倍角・三倍角の公式の右辺と同じ形をしている.
実際,
\[
  T_n(\cos\theta)=\cos(n\theta)
\]
が成り立つ.
導出を確かめよう.
$n=0,1$ では
\[
  T_0(\cos\theta)=1=\cos0,\qquad
  T_1(\cos\theta)=\cos\theta
\]
で一致する.
また余弦の加法定理から
\[
  \cos((n+1)\theta)
  =2\cos\theta\cos(n\theta)-\cos((n-1)\theta)
\]
である.
これは $T_{n+1}(\cos\theta)$ を作る漸化式と同じである.
したがって,二つの初めの値が一致し,次の値を作る規則も一致しているので,帰納法によりすべての $n$ で
\[
  T_n(\cos\theta)=\cos(n\theta)
\]
となる.
チェビシェフ多項式は,余弦の整数倍の角を,角度を使わずに $x=\cos\theta$ の多項式として記録している.

この関係は,本編で扱った三角関数型漸化式
\[
  a_{n+1}=2a_n^2-1
\]
にも戻ってくる.
初項が $-1\le a_1\le1$ なら,ある角 $\theta$ を選んで $a_1=\cos\theta$ と書ける.
二倍角の公式と $T_2(x)=2x^2-1$ から
\[
  a_{n+1}=T_2(a_n)
\]
であり,$\,a_n=\cos(2^{n-1}\theta)$ と予想できる.
この予想は初項で正しく,成り立つと仮定すれば
\[
  a_{n+1}
  =2\cos^2(2^{n-1}\theta)-1
  =\cos(2^n\theta)
\]
となるので,帰納法で確かめられる.
さらに $a_1=\cos\theta$ と $T_m(\cos\theta)=\cos(m\theta)$ を使えば
\[
  a_n=\cos(2^{n-1}\theta)
  =T_{2^{n-1}}(a_1)
\]
である.
たとえば本編の $a_1=\sqrt3/2=\cos(\pi/6)$ では
\[
  a_n=T_{2^{n-1}}\left(\frac{\sqrt3}{2}\right)
  =\cos\left(2^{n-1}\frac{\pi}{6}\right).
\]
つまり,本編の二次式による更新を何度も繰り返すことは,角度を二倍ずつ進めることと同じである.
その更新を多項式の言葉で書き直したものがチェビシェフ多項式だった.

より一般に,$\,T_m(T_k(x))=T_{mk}(x)$ という合成の関係も,同じ角度の式から分かる.
$x=\cos\theta$ とすれば
\[
  T_m(T_k(x))
  =T_m(\cos(k\theta))
  =\cos(mk\theta)
  =T_{mk}(x).
\]
この等式は $[-1,1]$ のすべての $x$ で成り立ち,両辺が多項式なので,多項式の恒等式として実数全体で成り立つ.
したがって $T_2$ を繰り返し合成する操作は,添字を二倍ずつ進める $T_2,T_4,T_8,\ldots$ にまとまる.

\subsection{非線形な更新も,不動点の近くでは線形に見える}\label{節：非線形な更新も,不動点の近くでは線形に見える}

ここまでは,漸化式を行列で表し,線形な更新の方向ごとの倍率を見てきた.
最後に,一次分数型のような非線形の更新規則でも,不動点の近くでは同じ見方ができることを確かめよう.
一項前から次が決まる漸化式
\[
  a_{n+1}=f(a_n)
\]
を考え,$\,\alpha$ を不動点 $f(\alpha)=\alpha$ とする.
そこからのずれを $e_n=a_n-\alpha$ と置くと
\[
  e_{n+1}=f(\alpha+e_n)-f(\alpha)
\]
である.
$f$ が $\alpha$ で微分可能なら,微分係数の定義から
\[
  \frac{f(\alpha+e)-f(\alpha)}{e}\longrightarrow f'(\alpha)
  \qquad(e\longrightarrow0)
\]
である.
つまり,ずれ $e$ が十分小さい範囲では
\[
  f(\alpha+e)-f(\alpha)\approx f'(\alpha)e
\]
となる.
曲線を不動点のすぐ近くまで拡大すると接線に見える,という一次近似を使ったのである.

この近似は,収束についても手がかりを与える.
もし $|f'(\alpha)|<1$ なら,$\,|f'(\alpha)|$ より大きく $1$ より小さい数 $c$ を選べる.
微分係数の定義により,$\,a_n$ が $\alpha$ に十分近ければ
\[
  |e_{n+1}|
  =|f(\alpha+e_n)-f(\alpha)|
  \le c|e_n|
\]
となる.
誤差は一歩ごとに少なくとも一定の割合で縮み,いったん十分近くに入ればその近くにとどまって $\alpha$ へ近づく.
したがって $|f'(\alpha)|<1$ は,不動点が近くの初期値を引き寄せることを示す.
先ほどの行列では固有値が各方向の倍率を表していた.
ここでは微分係数 $f'(\alpha)$ が,不動点近くでのずれの倍率を表している.
線形代数と微分は別々の話に見えるが,どちらも「一歩ごとの変化を倍率で測る」道具としてつながっている.
